% 图论（综述）
% license CCBYSA3
% type Wiki

本文根据 CC-BY-SA 协议转载翻译自维基百科\href{https://en.wikipedia.org/wiki/Graph_theory}{相关文章}

\begin{figure}[ht]
\centering
\includegraphics[width=6cm]{./figures/099d0167fdbd2081.png}
\caption{} \label{fig_TUl_1}
\end{figure}
在数学和计算机科学中，图论是对图的研究。图是一种数学结构，用于刻画对象之间的成对关系。在这一语境下，一个图由顶点（亦称为节点或点）组成，顶点之间通过边（亦称为弧、链接或线）连接。图论中区分无向图与有向图：在无向图中，边对两个顶点的连接是对称的；而在有向图中，边对两个顶点的连接是不对称的。图是离散数学研究的主要对象之一。
\subsection{定义}
图论中的定义不尽相同。以下是一些更基本的定义图及其相关数学结构的方式。
\subsubsection{图}
\begin{figure}[ht]
\centering
\includegraphics[width=6cm]{./figures/0674f0c8b9e7f86a.png}
\caption{一个具有三个顶点和三条边的无向图。} \label{fig_TUl_2}
\end{figure}
在一个受限但非常常见的意义下，[1][2]图是一个有序对$G = (V, E)$其中：
\begin{itemize}
\item $V$：顶点的集合（亦称为节点或点）；
\item $E \subseteq {{x, y} \mid x, y \in V ;\text{且}; x \neq y}$：边的集合（亦称为链接或线），其元素是无序的顶点对（即一条边关联两个不同的顶点）。
\end{itemize}
为了避免歧义，这类对象可称为无向单图。

在边${x, y}$中，顶点$x$与$y$称为该边的端点。该边被认为连接$x$与$y$，并与$x$、$y$关联。图中可能存在某些顶点没有属于任何边。根据该定义，不允许出现多重边（即两条或以上的边连接相同的一对顶点）。
\begin{figure}[ht]
\centering
\includegraphics[width=6cm]{./figures/e0f64ecda97f8953.png}
\caption{一个具有3个顶点、3条边和4条自环的无向多重图示例。} \label{fig_TUl_3}
\end{figure}
在另一种更一般的意义下，允许多重边，[3][4]图是一个有序三元组$G = (V, E, \phi)$其中：
\begin{itemize}
\item $V$：顶点的集合（也称为节点或点）；
\item $E$：边的集合（也称为链接或线）；
\item $\phi : E \to {{x,y} \mid x,y \in V ;\text{且}; x \neq y}$：一个关联函数，它将每一条边映射到一个无序的顶点对（即一条边连接两个不同的顶点）。
\end{itemize}
为了避免歧义，这类对象可称为无向多重图。

自环是一条将顶点连接到其自身的边。在前述两种定义下，图都不能有自环。因为一条将顶点$x$与自身连接的自环对应的边为${x,x} = {x}$，而它并不属于${{x,y}\mid x,y \in V, x \neq y}$。若要允许自环，则必须扩展定义：对于无向单图，其边集应修改为
$E \subseteq {{x,y}\mid x,y \in V}$；对于无向多重图，其关联函数应修改为
$\phi : E \to {{x,y}\mid x,y \in V}$。为了避免歧义，这些类型的对象可分别称为允许自环的无向单图和允许自环的无向多重图（有时也称为无向伪图 pseudograph）。

通常情况下，(V) 和 (E) 被认为是有限的；许多著名的结论在无限图中并不成立（或有所不同），因为许多推理在无限情形下失效。此外，通常假设 (V) 是非空的，但 (E) 可以是空集。

* 图的**阶（order）**是 (|V|)，即其顶点数。
* 图的**大小（size）**是 (|E|)，即其边数。
* 一个顶点的**度（degree 或 valency）**是与其关联的边的条数，其中自环计作两次。
* 图的**度**是所有顶点度数的最大值。

在一个阶为 (n) 的无向单图中，每个顶点的最大度为 (n-1)，该图的最大边数为

[
\frac{n(n-1)}{2}.
]

允许自环的无向单图 (G) 的边集，在其顶点集上诱导出一个**对称齐次关系** (\sim)，称为 (G) 的**邻接关系**。具体而言，对于每条边 ((x, y))，其端点 (x) 和 (y) 被称为**邻接**，记作

[
x \sim y.
]

---

要不要我帮你把“阶、大小、度、邻接”这些术语整理成一个简洁的术语表？
